\documentclass[12pt]{ximera}
\title{Homework 2: Sets, Subsets, and Set Operations}
\begin{document}
\begin{abstract}
Section 7.1: listing the elements of sets built from complements, differences, and
intersections; explaining a set identity in words; and counting the subsets of
a finite set.
\end{abstract}
\maketitle

\section*{Section 7.1: Sets, Subsets, and Set Operations}

\begin{problem}
Let
\[
\mathcal{U}=\set{2,3,4,5,7,9}, \quad
X=\set{2,3,4,5}, \quad
Y=\set{3,5,7,9}, \quad
Z=\set{2,4,5,7,9}.
\]
List the elements of each of the following sets. In each part, enter the elements in
increasing order.

\begin{prompt}
\textbf{(a)} $\olsi{X} = \set{\answer{7},\ \answer{9}}$

\begin{hint}
The complement $\olsi{X}$ is everything in the universal set $\mathcal{U}$ that is
\emph{not} in $X$.
\end{hint}

\begin{feedback}
Removing the elements of $X=\set{2,3,4,5}$ from $\mathcal{U}=\set{2,3,4,5,7,9}$ leaves
\[ \olsi{X} = \set{7,9}. \]
\end{feedback}
\end{prompt}

\begin{prompt}
\textbf{(b)} $Z - \olsi{X} = \set{\answer{2},\ \answer{4},\ \answer{5}}$

\begin{hint}
$Z-\olsi{X}$ means the elements of $Z$ with any element of $\olsi{X}$ thrown out.
\end{hint}

\begin{feedback}
$Z=\set{2,4,5,7,9}$ and $\olsi{X}=\set{7,9}$. Removing $7$ and $9$ from $Z$ leaves
\[ Z-\olsi{X} = \set{2,4,5}. \]
(Equivalently, $Z-\olsi{X}=Z\cap X$.)
\end{feedback}
\end{prompt}

\begin{prompt}
\textbf{(c)} $Y\cap\left(Z - \olsi{X}\right) = \set{\answer{5}}$

\begin{feedback}
From part (b), $Z-\olsi{X}=\set{2,4,5}$, and $Y=\set{3,5,7,9}$. The only element they
share is $5$, so $Y\cap\left(Z-\olsi{X}\right)=\set{5}$.
\end{feedback}
\end{prompt}
\end{problem}

\begin{problem}
Explain in words why
\[ \left(A\cap B\right)\cup\left(A\cap\olsi{B}\right) = A. \]

\begin{freeResponse}
\end{freeResponse}

\begin{hint}
Translate each piece into words using ``and'' and ``not''. Every element of $A$ either
is in $B$ or is not.
\end{hint}

\begin{feedback}
The left-hand side is the set of things that are ``in $A$ and in $B$'' together with
the things that are ``in $A$ and not in $B$''.

Every element of $A$ falls into exactly one of those two groups, since an element is
either in $B$ or it is not. So between them the two pieces contain all of $A$. And both
pieces are parts of $A$, so the union contains nothing outside $A$. Therefore the union
is exactly $A$.
\end{feedback}
\end{problem}

\begin{problem}
Suppose $6$ flavors of cat food are available at the store. Euclid the cat might like all $6$ flavors, or none of them, or any other
subset of them. How many possibilities are there for the subset of flavors Euclid
likes among the $6$ available flavors?

$\answer{64}$ possible subsets.

\begin{hint}
Decide flavor by flavor: for each of the $6$ flavors, Euclid either likes it or does
not. That is $2$ choices per flavor, made independently.
\end{hint}

\begin{feedback}
A set with $n$ elements has $2^n$ subsets, since each element is independently either
in or out. Here $n=6$, so there are
\[ 2^6 = 64 \]
possible subsets of flavors --- including the empty set (Euclid likes nothing) and
the whole set (Euclid likes everything).
\end{feedback}
\end{problem}

\begin{problem}
A concert features a cellist, a flutist, a harpist, and a vocalist. Each song
features a different subset of \emph{at least two} of the featured musicians. How
many songs are in the concert if every such subset is featured once?

$\answer{11}$ songs.

\begin{hint}
Start with all $2^4$ subsets of the four musicians, then subtract the ones that are
too small to be a song: the empty set, and the four solos.
\end{hint}

\begin{feedback}
There are $2^4=16$ subsets of the four musicians in total. The subsets that do
\emph{not} have at least two members are the empty subset ($1$ of them) and the
one-musician subsets ($4$ of them), for $5$ subsets to discard. That leaves
\[ 16 - 1 - 4 = 11 \]
subsets of at least two musicians, so the concert has $11$ songs.
\end{feedback}
\end{problem}

\end{document}
